\documentclass[12pt]{article}
\usepackage[utf8]{inputenc}
\usepackage[T1]{fontenc}
\usepackage{hyperref}
\title{Thudam Language Resources for K-12 Curriculum Development}
\author{Compiled by Claude Sonnet 5 (claude-sonnet-5) for Andr\'as Kornai}
\date{}
\begin{document}
\maketitle

\section{Introduction}
The code thw (Thudam) was retired from ISO 639-3 effective 2020-05-15 and merged into ola (Walungge). Older sources, such as the Rosetta Project, still list Thudam as a Tibetic variety of Nepal. The search found no separate language resources for Thudam, and no speaker numbers.

\section{Primary Educational Resources}
None found.

\section{Linguistic and Reference Materials}
Only code records and catalogue entries turned up, with no descriptive work: \url{https://rosettapanglossia.longnow.org/index.php?title=Thw} and \url{https://app.panlex.org/panlinx/ex/2739046}

\section{Religious and Community Resources}
None found.

\section{Digital and Audio Resources}
None found.

\section{Curriculum Development Context}
Thudam is now treated as part of Walungge (ola), so any resource search should use that code. The Thangmi (thf) results that appeared in the search are a different language and are not included here.

\end{document}